\section{Next work}

The first two steps below address the two sensitivities that moved the case's climate most, in order of
their size, and the third tests the reference both rest on; the rest are unordered, since their effects on the
climate remain unmeasured. Each names the question it settles, the comparison that settles it and the result
that would change the lunar design or the reach.

\textbf{The longwave at the settled state.} The largest single sensitivity rests on a correction fitted at
295~K and used near 307~K with twice the water vapour. Refitting the correction on the settled state's
cloud-free columns, and iterating to a consistent pair of fit and climate, decides whether the warm state
holds; a correction carrying humidity as well as temperature, and then a reduced correlated-$k$ scheme built
from the reference model's own spectroscopy in place of Sasamori's fits, would remove the fit altogether. If
the consistent state returns near 295~K, the design's dimming stands; if it stays near 307~K, the dimming that
holds the design temperature must grow, and the rain and cloud change with it.

\textbf{Cloud from the cloud-resolving model.} PlaSim's cloud amount and cloud water move the Moon by several
kelvin and have no lunar calibration. Forcing CM1 from the design case, with the transport across each ring
imposed, tests first whether the near-ground gap closes, and then gives cloud cover and water by pressure to
which PlaSim's humidity thresholds and its cloud-water profile can be fitted. Refitted cloud could move the
Moon as far as the cloud settings do, up to about 6~K, and the design's dimming with it. An independent GCM on
the same scenario, ROCKE-3D or the LMD generic model \autocite{way2017resolving,wordsworth2011gliese}, would
show which of the two present models departs near the ground: air as warm as CM1's would make PlaSim's
boundary layer and longwave the next repairs, and air as cool as PlaSim's would point to the slices' forcing
and CM1's microphysics.

\textbf{The reference against an independent code.} Every calibration rests on the line-by-line model, whose
carbon dioxide, oxygen, collision-induced absorption and sunlight no independent code has checked on the lunar
column. Running the design column through LBLRTM for the longwave and a DISORT-based solver for the sunlight,
with every absorber on identical layers, tests the reference itself. Layer heating or a split of sunlight
departing from the reference by more than the calibration residuals, 0.97~W\,m$^{-2}$ rms by layer or 0.6
points in the split, would require the multipliers and the gas bands to be refitted and the dimming reset.

\textbf{The tall air's dynamics.} Both dynamical cores hold gravity constant in a thin shell, where the lunar
air loses 5\% of its gravity by its mass-mean height and, in spherical shells, holds 11.5\% more air over each
square metre of ground than $p_s/g_s$. The consistent deep-atmosphere equations \autocite{white2005consistent}
and the deep forms of global models \autocite{mayne2014unified} include both effects; a first step corrects
the GCM's column mass and heights, and a test of the circulation's sensitivity shows whether the full
equations are needed. A change in the heat carried poleward larger than the spread between the present models
would make a deep core a requirement for the design case.

\textbf{Radiation in CM1.} RRTMG extrapolates its tables over the lowest 7.5~km of air over the sea, takes the
Sun's unfiltered spectrum, and cools the lowest 1.8~km 4.5~W\,m$^{-2}$ more than the reference. Extending its
tables past 1,053.63~hPa, or replacing it with the in-house correlated-$k$ scheme, and giving it the shield's
spectrum would separate the table's edge from the other causes. If extended tables remove most of the
4.5~W\,m$^{-2}$, the extrapolation was the cause and the near-ground range between the models narrows with it;
if they remove little, the cause lies in RRTMG's other approximations, and replacing RRTMG becomes the step.

\textbf{Microphysics at lunar gravity.} Morrison's limits on the size of rain, its breakup and the size of ice
crystals stay Earth's; snow falls with the drops' gravity exponent where its own mass and area laws give a
larger one; and in the boxes 85\% of the falling water evaporates, in the rings 57--73\%. A gravity-similarity
pair under prescribed radiative cooling would measure how far the scaled microphysics departs from
similarity, and limits scaled with the capillary length would show how much of the evaporation the small
drops cause.\claim{N03} If larger drops bring CM1's surface rain toward the GCM's, the rain gap closes on the
GCM's side and the rain-fed lakes of the design keep their source; if CM1's rain stays near a fifth of the
GCM's, the lakes' water budget has to be drawn from the cloud-resolving rain.

\textbf{A test of the reach.} The repairs have run on the Moon alone. Released and repaired
ExoPlaSim on the Earth-air slow rotators of Figure~\ref{fig:regimes}, with solar days of 16 and 64 days and the
released radiation, which suits their Earth-sized columns, would test whether the general repairs carry
without refitting. A general repair that needs retuning off the Moon would move to the calibrations, and the
reach of Table~\ref{tab:reach} would narrow with it.

\textbf{Electricity and seas.} The sign of the lunar lightning rests on laboratory charging laws that disagree
in the lunar storms' little cloud water. The sign of the charge rimed ice takes at 4--8~m\,s$^{-1}$ in
0.02--0.2~g\,m$^{-3}$ of cloud water at $-5$ to $-16$\,\textdegree C, conditions Earth's laboratories can
reach, would decide whether the lunar ground strikes are negative, as under Saunders and Peck's law, or
positive, as under Takahashi's, and so the polarity that protection on the ground and in the air must
meet.\claim{N01} A leader model would replace the bound on ground strikes, now between none and a fifth of
flashes, with a rate. The seas need the air's density field in place of one value, swell decay over the long
fetches where swell carries a third to two thirds of the wave energy,\claim{N04} a test of convergence in frequency for
the small seas that fail it, and wind input and breaking checked against theory at the lunar ratio of air to
water density; the last decides whether the similarity image of an Earth sea holds on the Moon at all.

\textbf{Reports and the record.} The Sun clock is reported to ExoPlaSim's maintainers and awaits a
reply.\claim{N02} The convective-cover unit, with its bias in synchronous runs, the orbit default and SWAN's
uninitialised breaker parameter await reports that would serve every user of those codes; CM1's thread race,
placed by elimination in its own loops, awaits locating and a report; and PlaSim's cold lowest level awaits a
test against smoothed terrain that would confirm or clear the spectral transform. CM1 would gain from one
gravity parameter honoured by every scheme. The record itself needs the twilight pair's control rerun on one
build, runtime identity runs for the Sun clock, the water scaling and patches 07--09 at their off settings,
and forcing for the CM1 cases set up before the Sun and the starting soundings were corrected.
